Man wende die Methode von Hermite auf die folgenden Integrale an
\begin{teilaufgaben}
\item $\displaystyle
\int \frac{2x^2+x+1}{(x+3)(x-1)^2}\,dx$
\item $\displaystyle
\int \frac{x^3+7x^2-5x+5}{(x-1)^2(x+1)^3}\,dx$
\end{teilaufgaben}

\begin{loesung}
\begin{teilaufgaben}
\item 
Zunächst ist der erweiterte euklidische Algorithmus auf die beiden
Nennerfaktoren anzuwenden:
\[
-\frac{1}{16}(x-5)(x+3)+\frac{1}{16}(x+1)^2 = 1.
\]
Multiplizieren mit dem Zähler gibt
\[
\begin{aligned}
s(x)
&=
-\frac{1}{16}(x-5)(x^2+x+1)
=
-\frac{1}{16}(
2 x^3 - 9 x^2  - 4 x - 5
)
\\
t(x)
&=
\frac{1}{16}(2x^2+x+1),
\end{aligned}
\]
die 
\[
s(x) (x+3) + t(x) (x-1)^2 = 2x^2 + x + 1
\]
erfüllen.
Division durch $(x+3)(x-1)^2$ ergibt
\[
\frac{s(x)}{(x-1)^2}
+
\frac{t(x)}{x+3}
=
\frac{2x^2+x+1}{(x+3)(x-1)^2}.
\]
Das gesuchte Integral ist daher
\begin{align*}
I
&=
\int \frac{s(x)}{(x-1)^2}\,dx
+
\int \frac{t(x)}{x+3}\,dx
\\
&=
-\int \frac{2x-5}{16}\,dx +\int \frac{x}{(x-1)^2}\,dx
+
\int \frac{2x-5}{x}\,dx + \int \frac{1}{x+3}\,dx
\end{align*}
Die beiden polynomiellen Integrale heben sich weg und das letzte
Integral hat bereits die gewünschte Form mit quadratfreiem Nenner.
Es ist also nur noch das zweite Integral mit der Methode von Hermite
zu reduzieren.
Dazu muss der grösste gemeinsame Teiler von $b(x)=x-1$ und $b'(x)=1$
in der Form
\[
s(x) b(x) + t(x) b'(x) = x
\]
dargestellt werden.
Die offensichliche Lösung ist $s(x)=1$ und $t(x)=1$.
Das Integral kann jetzt reduziert werden:
\begin{align*}
\int \frac{x}{(x-1)^2}\,dx
&=
\int\frac{1+1'/(2-1)}{x-1}\,dx
+
\frac{t(x)}{(1-2)(x-1)}
\\
&=
\int\frac{1}{x-1}\,dx
-
\frac{1}{x-1}
\\
&=
\log(x-1) - \frac{1}{x-1}.
\end{align*}
Damit kann man jetzt das gesuchte Integral zusammensetzen:
\begin{align*}
\int \frac{2x^2+x+1}{(x+3)(x-1)^2}\,dx
&=
-\frac{1}{x-1}+\log(x-1)
+\log(x+3).
\end{align*}

\item
Zunächst ist wieder die Zerlegung in Integrale mit Nennern
$(x-1)^2$ und $(x+1)^3$ vorzunehmen.
Dazu verwendet man
\[
\frac{1}{16}(3x^2+10x+11)(x-1)^2
-
\frac{1}{16}(3x-5)(x+1)^3
=
1
\]
Es folgt, dass der Zähler $x^3+7x^2-5x+5$ des Integranden
\[
\frac{1}{16}(3x^5+31x^4+33x^3+21x^2-5x+55)(x-1)^2
-
\frac{1}{16}(3x^4+16x^3-50x^2+40x-25)(x+1)^3
%=
%x^3+7x^2-5x+5
\]
ist.
Im gesuchten Integral heben sich die polynomiellen Teile wieder 
weg und es bleibt
\begin{align*}
\int \frac{x^3+7x^2-5x+5}{(x-1)^2(x+1)^3}\,dx
&=
\int\frac{4}{(x+1)^3}\,dx
+
\int\frac{1}{(x-1)^2}\,dx
\end{align*}
Auf beide Integrale muss jetzt noch der Satz von Hermite
angewendet werden.

Für das zweite Integral ist $b(x)=x-1$ und $b'(x)=1$ mit
\[
s(x)b(x)+t(x)b'(x)=1
\qquad\Rightarrow\qquad
=
s(x)=0, t(x)=1
\]
und daher
\[
\int \frac{1}{(x-1)^2}\,dx
=
-
\frac{1}{x-1}.
\]

Für das erste Integral ist $b(x)=x+1$ und $b'(x)=1$ mit der gleichen
Lösung $s(x)=0$ und $t(x)=4$.
Nach dem Satz von Hermite ist dann
\begin{align*}
\int \frac{4}{(x+1)^3}\,dx
&=
\frac{4}{(1-3)(x+1)^2}
=
-
\frac{2}{(x+1)^2}.
\end{align*}
Insgesamt finden wir daher für das gesuchte Integral
\[
\int \frac{x^3+7x^2-5x+5}{(x-1)^2(x+1)^3}\,dx
=
-
\frac{1}{x-1}
-
\frac{2}{(x+1)^2}
=
-
\frac{x^2+4x+1}{x^3+x^2-x-1}.
\qedhere
\]
\end{teilaufgaben}
\end{loesung}
